%!TeX root=../monografia
%("dica" para o editor de texto: este arquivo é parte de um documento maior)
% para saber mais: https://tex.stackexchange.com/q/78101
\chapter{Introdução}
\label{cap:introducao}

\section{Introdução}

\subsection{Problema do predecessor}

O problema do predecessor tem como objetivo manter um conjunto $S$ de $n$
elementos de um universo $U$ sujeito às seguintes operações:
\[
\Predecessor(x) = \max \{y \in S \mid y < x \} 
\]
\[
\Sucessor(x) = \min \{y \in S \mid y > x \}
\]
%(Mauricio): x \in U nas consultas também? Imagino que sim.

No nosso caso, estamos interessados no problema dinâmico, que também suporta as
operações de \textbf{Inserção}, para $x \in U$, e \textbf{Remoção}, para $x \in
S$.

Com uma Árvore Binária de Busca Balanceada (ABBB), conseguimos resolver o
problema com custo $\Theta (\lg n)$ para as operações e ocupando espaço $\Theta
(n)$.
%(Mauricio): Aqui você já está supondo Word RAM, visto que está considerando que os ponteiros de uma AABB têm tamanho O(1) cada.
No modelo clássico baseado em comparações, as únicas operações entre os valores
dos elementos são de comparação (como $\geq, =, \leq$) e nele uma ABBB é
assintoticamente ótima.

\subsection{Modelo Word RAM}

Vamos usar o modelo \textbf{word RAM} (Random Access Machine), definido por
Fredman e Willard em 1983 [conferir veracidade]. Nele, a memória é dividida em
\textit{palavras} com tamanho de $w$ bits. Assim, todos os elementos são
inteiros de $w$ bits e definimos um universo finito $U=\{0,1,\dots,u-1\}$, com
$u=2^w$.  

As operações aritméticas e lógicas entre elementos são executadas em tempo
constante:
\[+, -, *, /,\%,<<,>>,\&,|,\geq,=,\leq\]

Além disso, temos que $w\geq \lg n$, caso contrário, nem seria possível
representar o problema nesse modelo. [desenvolver]
%(Mauricio): Acho que $n$ deveria ser $u$.

Também vale notar que esse modelo já se aproxima da arquitetura real dos
computadores hoje em dia. Atualmente, um sistema comum tem $64$ bits, ou seja,
suportaria um universo de palavras de $0$ a $2^{64}$.
